\documentclass[11pt]{article}
\setlength{\textwidth}{160mm}
\setlength{\textheight}{235mm}
\setlength{\oddsidemargin}{0mm}
\setlength{\topmargin}{-12mm}
\usepackage{amsmath,amssymb,amsthm}
\usepackage[T1]{fontenc}
\usepackage[colorlinks=true,linkcolor=blue,urlcolor=blue]{hyperref}
\newtheorem{theorem}{Theorem}
\newtheorem{lemma}{Lemma}
\newtheorem{corollary}{Corollary}
\newtheorem{remark}{Remark}
\newcommand{\R}{\mathbb R}
\newcommand{\E}{\mathbb E}
\newcommand{\eps}{\varepsilon}
\newcommand{\aff}{\operatorname{aff}}
\newcommand{\dist}{\operatorname{dist}}
\newcommand{\Cov}{\operatorname{Cov}}
\newcommand{\diver}{\operatorname{div}}
\newcommand{\tr}{\operatorname{tr}}
\title{Extensions of empirical-flow atomicity:\\ schedules, perturbations, and a failure under composition}
\author{Research proof supplement}
\date{17 October 2026}
\begin{document}
\maketitle
\begin{abstract}
This supplement proposes extensions of the two atomicity propositions in the
current manuscript. A quantitative hypothesis allows time-dependent guidance
and vanishing scaled field errors while preserving almost-sure atomicity.
A separate exact construction shows that averaging conditional fields and then
applying guidance can instead create atoms outside the training set. Its scalar
threshold is sharp, its off-training basin has positive measure, and a tensor
construction has full affine span. The proof arguments have been independently
reviewed; their priority and usefulness for a conference
submission require a separate literature and experimental assessment.
\end{abstract}

\section{Common finite-support fields}
Let $A=\{a_1,\ldots,a_N\}\subset\R^d$ have distinct elements. For $s>0$ put
\[
 t=1-e^{-s},\qquad \eps=e^{-s},\qquad
 \alpha=\eps^2/t,\qquad \beta=\alpha^{-1}=e^{2s}-e^s,
 \qquad b_j=x\cdot a_j-\tfrac t2|a_j|^2.
\]
For a strictly positive probability vector $p$ define
\[
 q_{p,j}=\frac{p_j e^{\beta b_j}}{\sum_kp_ke^{\beta b_k}},\qquad
 m_p=\sum_jq_{p,j}a_j,\qquad
 H_p=\alpha\log\sum_jp_je^{b_j/\alpha}.
\]
The probabilities and means extend smoothly to $s=0$ by $\beta(0)=0$.
Fix two such vectors $p_c,p_u$. For a nonnegative guidance schedule $w$ let
\[
 m_w=(1+w)m_c-wm_u,\qquad f(x,s)=m_w-x.
\]
This is the logarithmic-time version of velocity $(m_w-x)/(1-t)$.
For $z\in\R^d$ let $N(z)$ be its nearest atoms and set
\[
 C_A=\{z:z\in\aff N(z)\}.
\]
This set is finite: for each nonempty subset $S\subset A$, two points in
$\aff S$ equidistant from every member of $S$ have difference both parallel
and perpendicular to $\operatorname{span}(S-S)$, hence coincide.

\clearpage
\section{A sufficient robustness theorem}
\begin{theorem}[Scheduled guidance with scaled perturbations]\label{thm:robust}
Assume $w:[0,\infty)\to[0,W]$ is locally absolutely continuous and, for some
$s_0>0$,
\[
 \int_{s_0}^{\infty}\alpha(s)|w'(s)|\,ds<\infty.
\]
Let $R(x,s)$ be continuous in $(x,s)$, continuously differentiable in $x$,
with spatial derivative locally bounded on compact subsets. Suppose
\[
 |R(x,s)|\le\rho(s),\quad \rho\text{ locally bounded},\quad
 \rho(s)\to0,\quad \int_{s_0}^\infty\rho(s)^2\,ds<\infty,
\]
uniformly in $x$, and suppose there is a locally bounded nonnegative $\eta$
with $\eta(s)\to0$ such that, for $s\ge s_0$,
\[
 \diver R(x,s)\ge-\eta(s)\beta(s)\qquad\text{for every }x.
\]
Every trajectory of $x'=f(x,s)+R(x,s)$ from a finite point at $s=0$ exists
globally and converges to a point of $C_A$. The set of initial points whose
limits are outside $A$ is Lebesgue-null. Every absolutely continuous source
therefore has an endpoint law supported on $A$.
\end{theorem}
\begin{proof}
Write $M=\max_j|a_j|$. Then $|m_w|\le(1+2W)M$ and variation of constants in
$x'=-x+m_w+R$ bounds every trajectory, and all trajectories from a bounded
initial set, uniformly in future time. The locally bounded envelope also gives
linear growth on each finite interval, forward and backward. Local spatial
regularity gives unique global solutions and finite-time $C^1$ diffeomorphisms.

Put $U=|x|^2/2-(1+w)H_c+wH_u$, so $\nabla U=-f$. For fixed $x$,
\[
 \partial_s H_p=-\alpha'D(q_p\Vert p)
 -\frac{e^{-s}}2\sum_jq_{p,j}|a_j|^2,
 \qquad 0\le D(q_p\Vert p)\le\log(1/\min_jp_j).
\]
Comparison of partition functions gives
$|H_c-H_u|\le L\alpha$, where
$L=\max_j|\log(p_{c,j}/p_{u,j})|$. Since $\alpha$ decreases to zero,
$|\alpha'|$ is integrable on $[s_0,\infty)$. It follows that
\[
 |\partial_s U|\le K(|\alpha'|+e^{-s})+L\alpha|w'|=:h(s),
 \qquad h\in L^1(s_0,\infty).
\]
The chain rule holds almost everywhere. Along the perturbed flow,
\[
 \frac d{ds}U(x(s),s)
 \le h(s)-\tfrac12|f(x(s),s)|^2+\tfrac12\rho(s)^2.
\]
The min--max bounds $\min_jb_j\le H_p\le\max_jb_j$ show that $U$ is bounded
below on each bounded trajectory, uniformly in $s$. Thus
\begin{equation}\label{eq:energy}
 \int_{s_0}^\infty|f(x(s),s)|^2\,ds<\infty.
\end{equation}

If $z\notin C_A$, its distance to $\aff N(z)$ is positive. In a small fixed
neighbourhood of $z$, every atom outside $N(z)$ has a positive limiting score
gap from at least one member of $N(z)$. Fixed positive priors make its softmax
mass tend uniformly to zero there. Bounded $w$ implies that $m_w$ approaches
$\aff N(z)$, so $|f|$ is bounded below on a smaller neighbourhood at late
times. Actual speed is eventually bounded because $x,m_w,\rho$ are bounded.
If $z$ were a cluster point, repeated visits to the inner half of that
neighbourhood would either remain there forever or supply disjoint residence
intervals of a fixed positive duration. Both alternatives contradict
\eqref{eq:energy}. Hence all cluster points lie in $C_A$. The omega-limit set
is nonempty, compact and connected, being the intersection of the closures of
bounded continuous trajectory tails. Since $C_A$ is finite, the trajectory
converges to one point of $C_A$.

Suppose this point is $z\notin A$, and put $r=\dist(z,A)>0$. Write
$u=m_c-z$, $v=m_u-z$, $\mu=m_w-z$ and
$T=(1+w)\tr\Cov_c(a)-w\tr\Cov_u(a)$. Concentration on $N(z)$ gives,
along the entire trajectory,
\[
 T=r^2-|\mu|^2+w(1+w)|u-v|^2+o(1)
 \ge r^2-|\mu|^2+o(1).
\]
The error tends to zero even for oscillating $w$, since $w$ is bounded.
Spatial differentiation gives $\diver f=\beta T-d$. Since
$\mu=f+x-z$, the divergence hypothesis implies
\[
 \diver(f+R)(x(s),s)
 \ge\beta(s)\{r^2-2|f(x(s),s)|^2-\zeta(s)\}-d,
 \qquad\zeta(s)\to0.
\]
We use three elementary weighted-tail estimates:
\[
 \int_{s_0}^s\beta=\tfrac12\beta(s)+o(\beta(s)),\quad
 \int_{s_0}^s\beta|f|^2=o(\beta(s)),\quad
 \int_{s_0}^s\beta\zeta=o(\beta(s)).
\]
The first follows by integrating $e^{2s}-e^s$. For the second, split at $S$;
the initial part divided by $\beta(s)$ tends to zero, and the tail is at most
$\beta(s)\int_S^\infty|f|^2$. Then send $S\to\infty$ using
\eqref{eq:energy}. For the third, split similarly and use
$\sup_{u\ge S}|\zeta(u)|\to0$ and the first estimate. Consequently, the
Jacobian $J_s$ of the finite-time flow initialized at $s_0$ satisfies, at each such source point,
\[
 \log J_s=\int_{s_0}^s\diver(f+R)(x(u),u)\,du
 \ge\tfrac12r^2\beta(s)+o(\beta(s))-O(s)\longrightarrow\infty.
\]
In the following measure argument source sets are at time $s_0$.
All trajectories converge, so their endpoint map is Borel. On the measurable
non-atom basin $E$ intersected with a bounded initial ball $B$, change of
variables and the common bounded image ball $K_B$ give
$\int_{E\cap B}J_s\le\operatorname{Leb}(K_B)$. Fatou along integer times
forces $\operatorname{Leb}(E\cap B)=0$. Exhaust the space by balls. The
finite-time diffeomorphism from $0$ to $s_0$ has a locally Lipschitz inverse,
which transfers this nullity to time zero. Absolute continuity gives the
claimed support statement, and bounded convergence for bounded continuous
test functions identifies the weak endpoint law.
\end{proof}

\begin{corollary}[Concrete schedules and errors]\label{cor:examples}
The theorem includes $R=0$ and schedules that do not converge, for example
$w(s)=1+\frac12\sin(e^s)$. It also includes
\[
 |R(x,s)|+\|D_xR(x,s)\|\le C e^{-\gamma s},\qquad \gamma>0.
\]
If an original-time field error is $\Delta$, then $R=(1-t)\Delta$.
Thus the sufficient condition
$|\Delta|+\|D_x\Delta\|\le C(1-t)^{-1+\gamma}$ permits original-time
errors that diverge when $0<\gamma<1$.
If a trajectory reaches an atom $a$ under this exponential envelope, then
\[
 |x(s)-a|=
 \begin{cases}
 O(e^{-s}),&\gamma>1,\\
 O(se^{-s}),&\gamma=1,\\
 O(e^{-\gamma s}),&0<\gamma<1.
 \end{cases}
\]
\end{corollary}
\begin{proof}
The schedule derivative has magnitude at most $e^s/2$ and
$\alpha e^s=1/(e^s-1)$ is integrable away from zero. For exponential errors,
the amplitude is square integrable and tends to zero, while
$\diver R\ge-dCe^{-\gamma s}=-o(\beta)$.
Near a reached isolated atom all competing logits have a uniform positive
gap, so fixed positive priors and bounded $w$ give
$|m_w-a|\le K e^{-c\beta(s)}$ for large $s$. In $y=x-a$ the equation is
$y'=-y+(m_w-a)+R$. Variation of constants and
$\int_S^\infty e^u e^{-c\beta(u)}du<\infty$ give the three rates directly.
\end{proof}

\begin{remark}[Scope of the certificate]
The hypotheses are uniform field controls, not consequences of a measured
teacher loss. In particular, a fixed neural field that remains regular up to
$t=1$ need not approximate the singular exact field in this scaled norm.
The theorem does not assert equality of endpoint probabilities or posterior
calibration. Finite observations of $R$ and its divergence do not certify its
uniform assumptions. Finitely many jumps of the schedule before an
absolutely continuous tail can be treated by applying the theorem after the
last jump and transferring the source through the finite-time flow.
\end{remark}

\section{A sharp failure when averaging is followed by guidance}
Take scalar atoms and labels $-1,+1$ with uniform empirical index, independent
source $N(0,1)$, and observed label $Y=a+hZ$, where $Z\sim N(0,1)$ is
independent and $h>0$. Put $b=h^{-2}$.
Fix a constant guidance strength $w\ge0$.
The exact label-conditioned posterior means at queries $\pm1$ are
$\tanh(\beta x\pm b)$; the unconditional mean is $\tanh(\beta x)$.
Apply CFG to the average of the two conditional fields:
\[
 v_{\mathrm{comp}}=\frac{1+w}{2}(v_++v_-)-wv_u.
\]
Its logarithmic-time dynamics is $x'=F(\beta x)-x$, where
\begin{equation}\label{eq:F}
 F(u)=\frac{1+w}{2}\{\tanh(u+b)+\tanh(u-b)\}-w\tanh u.
\end{equation}
This is guidance of an averaged field, which is generally not the exact field
for a single positive prior. The robustness theorem does not identify it with
such a field.

\begin{theorem}[Off-training mass has a sharp threshold]\label{thm:phase}
Let $w_c=\sinh(b)^{-2}$. Every trajectory of
$x'=F(\beta x)-x$ exists for all $s\ge0$ and has an endpoint.
If $0\le w\le w_c$, the endpoint is $\operatorname{sgn}(x(0))$,
with $\operatorname{sgn}(0)=0$.
If $w>w_c$, there is $R_*\in(0,\infty)$ such that the endpoint is
$-1$ on $(-\infty,-R_*)$, $0$ on $[-R_*,R_*]$, and $1$ on
$(R_*,\infty)$. For a standard Gaussian source and its cdf $\Phi$, the
supercritical endpoint law is
\[
 (1-\Phi(R_*))(\delta_{-1}+\delta_1)
 +(2\Phi(R_*)-1)\delta_0.
\]
Thus positive averaging and individual single-prior CFG preserve training
support, whereas their specified composition can put positive mass outside it.
\end{theorem}
\begin{proof}
Put $M=1+2w$. Boundedness $|F|\le M$ gives
$|x(s)|\le e^{-s}|x(0)|+M(1-e^{-s})$, hence global existence by local
smoothness and continuation. Oddness makes $0$ invariant; uniqueness preserves
the sign of a nonzero source. For $s>0$ let $u=\beta x$. Direct differentiation
gives
\[
 u'=c(s)u+\beta(s)F(u),\qquad
 c(s)=\frac{e^s}{e^s-1}>1,\qquad
 \frac{\beta(s)}{c(s)}=(e^s-1)^2.
\]
With $q=\tanh b$ and $z=\tanh u$, the addition formula gives
\begin{equation}\label{eq:factor}
 F(u)=\frac{z\{1-(1+w)q^2+wq^2z^2\}}{1-q^2z^2}.
\end{equation}
For $w\le w_c$ this is strictly positive when $u>0$, including equality
$w=w_c$, where its numerator is $wq^2z^3>0$. From any $s_1>0$, a positive
source has $u'\ge u$, hence $u\to\infty$ and $F(u)\to1$. Variation of
constants in $x'=F(u)-x$ gives $x\to1$. Oddness deals with negative sources.

Suppose $w>w_c$. Set
\[
 \kappa=\frac{w\sinh^2b-1}{\cosh^2b}>0,\quad
 \delta=\frac\kappa{4M},\quad
 S=\log(1+2/\sqrt\kappa),\quad \beta_S=\beta(S).
\]
We have $F'(0)=-\kappa$ and $|F''|\le2M$, so
$uF(u)\le-\kappa u^2/2$ for $|u|\le\delta$. For all $s\ge S$,
$\kappa\beta/2\ge2c$. At either boundary $|u|=\delta$,
\[
 \tfrac12(u^2)'\le(c-\kappa\beta/2)\delta^2<0.
\]
Thus $|u|\le\delta$, or $|x|\le\delta/\beta(s)$, is a forward-invariant
shrinking tube after $S$. Its trajectories have endpoint zero.
Since $|F'|\le M$, continuous dependence can be quantified by
\[
 |x(S)|\le |x(0)|e^{H},\quad
 H=S+M\left\{\frac{e^{2S}-1}{2}-(e^S-1)\right\}.
\]
Consequently every source with
$|x(0)|<\delta e^{-H}/(2\beta_S)=:\eps_*>0$
enters the tube.

Formula~\eqref{eq:factor} has precisely one positive nonzero root,
\[
 r_*=\operatorname{arcosh}(\sqrt w\sinh b),
\]
with $F<0$ on $(0,r_*)$ and $F>0$ on $(r_*,\infty)$.
For a positive source, either $u<r_*$ at every positive time, in which
case $0<x<r_*/\beta\to0$, or it reaches $r_*$. At that boundary
$u'=cr_*>0$, so $u\ge r_*$ is forward invariant and $u'\ge u$ there.
It follows that $u\to\infty$ and $x\to1$. All positive sources therefore
converge to $0$ or $1$, and oddness proves convergence for all sources.

Scalar uniqueness preserves order, so the endpoint map is nondecreasing.
Its zero set is an interval and is symmetric. The endpoint $1$ is equivalent
to $u>r_*$ at some finite time, an open condition on the source by continuous
dependence; the same holds for endpoint $-1$. Hence the zero basin is closed.
It contains $(-\eps_*,\eps_*)$. Finally,
$x(S)\ge e^{-S}x(0)-M(1-e^{-S})$ shows that sources larger than
$e^S r_*/\beta_S+M(e^S-1)$ escape to $1$. The zero basin is therefore
$[-R_*,R_*]$ with $0<R_*<\infty$. Pushing forward the Gaussian source gives
the displayed law.
\end{proof}

\begin{lemma}[The off-training basin does not require exact symmetry]\label{lem:tube}
Let $G:\R\to\R$ be bounded and locally Lipschitz. If there exist $l<r$
with $G(l)>0$ and $G(r)<0$, then a nonempty open interval of sources for
$x'=G(\beta(s)x)-x$ converges to zero. In particular this conclusion for
the supercritical $F$ persists under sufficiently small bounded locally
Lipschitz perturbations of $F$, even when they are not odd.
\end{lemma}
\begin{proof}
Bounded $G$ gives finite-time existence forward and backward. Uniqueness and
continuous dependence make each finite-time flow an increasing homeomorphism
of $\R$. In $u=\beta x$ coordinates, $u'=cu+\beta G(u)$. Since $c$ is
bounded at late times and $\beta\to\infty$, this field is strictly positive
at $l$ and strictly negative at $r$ for all sufficiently late $s$. Thus
$[l,r]$ is invariant after some $S$. The nonempty interval
$(l/\beta(S),r/\beta(S))$ has a nonempty open inverse under the finite-time
flow, and all its sources have endpoint zero. For supercritical $F$ the
previous $\delta$ gives $F(-\delta)\ge\kappa\delta/2$ and
$F(\delta)\le-\kappa\delta/2$. Uniform perturbations smaller than
$\kappa\delta/2$ preserve these two signs. For a source charging the open
inverse interval, including a Gaussian, this gives a positive-mass endpoint
submeasure. No assertion about convergence of other trajectories is needed
for that conclusion.
\end{proof}

\clearpage
\begin{corollary}[A full-dimensional empirical counterexample]\label{cor:tensor}
Let $A=\{-1,+1\}^d$ with uniform weights, observed label $y(a)=a_1$,
Gaussian label noise of standard deviation $h>0$, and independent source
$N(0,I_d)$. Average the exact fields at queries $+1,-1$ and apply CFG against
the unconditional exact field as in~\eqref{eq:F}. For $w>w_c$, the endpoint
law is
\[
 \mu_1\otimes(\tfrac12\delta_{-1}+\tfrac12\delta_1)^{\otimes(d-1)},\quad
 \mu_1=(1-\Phi(R_*))(\delta_{-1}+\delta_1)
 +(2\Phi(R_*)-1)\delta_0.
\]
It assigns mass $2\Phi(R_*)-1>0$ to the off-training set
$\{0\}\times\{-1,+1\}^{d-1}$. The training support has full affine span;
$N=2^d>d$ for $d\ge2$.
\end{corollary}
\begin{proof}
For a query $y=\pm1$, posterior weights are proportional to
\[
 \exp\left(-\frac{|x-ta|^2}{2\eps^2}
           -\frac{(y-a_1)^2}{2h^2}\right)
 \propto\exp\left((\beta x_1+yb)a_1+\sum_{k=2}^d\beta x_ka_k\right).
\]
They factor over coordinates, giving mean
$(\tanh(\beta x_1+yb),\tanh(\beta x_2),\ldots,\tanh(\beta x_d))$.
The unconditional mean removes $yb$. Hence the first coordinate follows the
scalar composed equation; every other coordinate follows
$x_k'=\tanh(\beta x_k)-x_k$. Its endpoint is the source sign by the
subcritical argument above. Independence of Gaussian source coordinates
proves the product law. The hypercube has full affine span, and no point with
first coordinate zero is a training atom.
\end{proof}

\section{Interpretation and next evidence}
The first theorem is a sufficient stability result using log-partition energy,
finite nearest-atom geometry, signed covariance and a Jacobian-volume
argument. Its assumptions must not be inferred from finite teacher loss.
The sharp composition theorem uses scalar comparison and an invariant
shrinking tube; the tube lemma explains why the failure is stronger than a
single symmetric exceptional source. The result creates new \emph{atoms},
not an absolutely continuous posterior and not a generalisation guarantee.

These extensions are kept separate from the existing manuscript while
their literature delta and experimental relevance are assessed. Useful tests are
the predicted threshold $w_c(h)=\sinh(h^{-2})^{-2}$, the positive zero-basin
mass, tensor-support controls, asymmetric priors, and resolution/stiffness
checks. A separate scientific note records literature proximity and which
parts remain hypotheses. No numerical experiment is used as a proof here.
\end{document}
